\documentclass[12pt]{article}
\usepackage[utf8]{inputenc}
\usepackage[T1]{fontenc}
\usepackage[margin=1in]{geometry}
\usepackage{setspace}
\usepackage{xcolor}
\usepackage[colorlinks=true, urlcolor=blue!60!black, linkcolor=blue!60!black]{hyperref}
\setlength{\parindent}{2em}
\setlength{\parskip}{1.25em}
\renewcommand{\baselinestretch}{1.0}
\title{\textit{Ex parte McCardle}, 74 U.S. (7 Wall.) 506 (1869)}
\author{Case Brief}
\date{}

\begin{document}
\maketitle

\section{Case Metadata}
\begin{description}
  \item[Citation:] \textit{Ex parte McCardle}, 74 U.S. (7 Wall.) 506 (1869) (December Term, 1868).
  \item[Argued:] March 2--5 \& 9, 1868 (merits); March 2 \& 19, 1869 (effect of the repealing act).
  \item[Decided:] April 12, 1869.
  \item[Vote:] 8--0 (unanimous; eight Justices sat on the Court: Chase, C.J., and Nelson, Grier, Clifford, Swayne, Miller, Davis, and Field, JJ.).
  \item[Author:] Chief Justice Salmon P. Chase.
  \item[Concurrences / Dissents:] None.
  \item[Chapter / Topic:] Judicial Review and the Federal Courts --- Congressional Control of Federal Appellate Jurisdiction (Exceptions Clause, U.S. Const.\ art.\ III, \S~2, cl.~2).
  \item[Primary Source URLs:] \href{https://supreme.justia.com/cases/federal/us/74/506/}{Justia U.S.\ Supreme Court Center} (\url{https://supreme.justia.com/cases/federal/us/74/506/}); \href{https://www.law.cornell.edu/supremecourt/text/74/506}{Cornell LII} (\url{https://www.law.cornell.edu/supremecourt/text/74/506}).
\end{description}

\section{Facts and Procedural History}
Following the Civil War, Congress enacted the Military Reconstruction Acts of 1867, placing the former Confederate states under military government. William H. McCardle, a civilian newspaper editor in Vicksburg, Mississippi, published a series of sharply critical editorials denouncing Reconstruction. Federal military authorities arrested McCardle and held him in military custody for trial before a military commission on charges arising from the publication of allegedly incendiary and libelous articles.

McCardle petitioned the U.S.\ Circuit Court for the Southern District of Mississippi for a writ of habeas corpus under the Habeas Corpus Act of February 5, 1867, which authorized federal courts to grant writs of habeas corpus whenever any person was restrained of liberty in violation of the Constitution, treaties, or laws of the United States, and expressly authorized appeals from Circuit Court habeas judgments to the U.S.\ Supreme Court. After a hearing, the Circuit Court remanded McCardle to military custody. McCardle appealed to the Supreme Court under the 1867 Act and was admitted to bail pending review.

At the prior term, the Supreme Court denied the government's motion to dismiss the appeal, confirming its jurisdiction under the 1867 Act (\textit{Ex parte McCardle}, 73 U.S. (6 Wall.) 318 (1868)). The Court then heard four days of full oral argument on the merits in early March 1868 and took the case under advisement. Fearing that the Court would invalidate the Military Reconstruction Acts (following \textit{Ex parte Milligan}, 71 U.S. (4 Wall.) 2 (1866)), Congress swiftly passed an amendment---enacted on March 27, 1868 over President Andrew Johnson's veto (15 Stat.\ 44)---repealing so much of the Act of February 5, 1867 ``as authorized an appeal from the judgment of the Circuit Court to the Supreme Court of the United States, or the exercise of any such jurisdiction by said Supreme Court, on appeals which have been, or may hereafter be taken.'' The Court continued the case while Chief Justice Chase presided over President Johnson's impeachment trial and subsequently heard reargument solely on the constitutional effect of the 1868 repealing statute.

\section{Issue(s) Presented}
\begin{enumerate}
  \item Does Article III, Section 2, Clause 2 of the U.S.\ Constitution---the Exceptions and Regulations Clause---empower Congress to strip the Supreme Court's statutory appellate jurisdiction over a habeas corpus appeal while that appeal is already pending and has been argued on the merits?
  \item Does the repealing Act of March 27, 1868 eliminate all Supreme Court appellate jurisdiction in habeas corpus cases, or only appeals brought under the Act of February 5, 1867?
\end{enumerate}

\section{Holding}
Yes to the first question; the repealing Act eliminates only appeals under the 1867 Act. Appeal \textbf{dismissed for want of jurisdiction}. Under the Exceptions Clause of Article III, Section 2, Congress possessed express constitutional authority to repeal the 1867 grant of appellate jurisdiction even while McCardle's appeal was pending after argument, and the Court may not inquire into Congress's legislative motives. Once statutory jurisdiction is repealed, the Court lacks authority to pronounce judgment and must dismiss the appeal. However, the 1868 Act did not disturb the Supreme Court's preexisting appellate habeas jurisdiction under the Judiciary Act of 1789.

\section{Reasoning}
Chief Justice Chase's concise unanimous opinion proceeds in four analytical steps:
\begin{enumerate}
  \item \textbf{Constitutional Source Subject to Express Congressional Exceptions:} Chase acknowledged that the Supreme Court's appellate jurisdiction is conferred directly by Article III of the Constitution rather than derived from statutes. Yet the constitutional grant is expressly qualified: it is conferred ``with such Exceptions, and under such Regulations as the Congress shall make.'' Relying on \textit{Durousseau v. United States}, 10 U.S. (6 Cranch) 307, 312 (1810), and \textit{Wiscart v. Dauchy}, 3 U.S. (3 Dall.) 321 (1796), Chase explained that Congress's affirmative statutory descriptions of appellate jurisdiction implicitly negate all appellate jurisdiction not described. Here, the exception was not even implied: the 1868 Act repealed the 1867 appellate provision in express terms---``hardly possible to imagine a plainer instance of positive exception.''
  \item \textbf{Irrelevance of Legislative Motive:} Addressing McCardle's contention that Congress acted specifically to shield Reconstruction from constitutional review, Chase declared that the Court is ``not at liberty to inquire into the motives of the legislature.'' Judicial review extends solely to constitutional power, and the power to make exceptions to appellate jurisdiction is granted by express constitutional text.
  \item \textbf{Effect of Repeal on Pending Appeals:} Rejecting petitioner's argument that stripping jurisdiction mid-appeal constituted an unconstitutional legislative exercise of judicial power, Chase distinguished cases in which legislatures interfered with courts exercising continuing jurisdiction. Citing \textit{Norris v. Crocker}, 54 U.S. (13 How.) 429 (1851), and \textit{Insurance Co. v. Ritchie}, 72 U.S. (5 Wall.) 541 (1867), Chase applied the settled rule that when a jurisdictional statute is repealed, it must be treated---except as to past and closed transactions---as if it had never existed. Without jurisdiction, the Court has no power to declare the law and can only announce that fact and dismiss.
  \item \textbf{Preservation of Alternative Preexisting Habeas Jurisdiction:} In the opinion's pivotal final paragraph, Chase rejected counsel's assumption that upholding the 1868 Act destroyed the Court's entire appellate power in habeas corpus cases. The 1868 Act repealed only appeals from Circuit Courts under the 1867 Act and did ``not affect the jurisdiction which was previously exercised''---leaving intact the Court's authority to issue ``original'' writs of habeas corpus and certiorari under Section 14 of the Judiciary Act of 1789.
\end{enumerate}

\section{Concurring and Dissenting Opinions}
None. All eight sitting Justices joined Chief Justice Chase's opinion. (Historically, when the Court postponed decision in April 1868 after Congress passed the repealing statute, Justice Grier---joined by Justice Field---entered a sharp protest against the delay, warning that the Court was evading its duty to decide a case already argued; nevertheless, both Justices joined the unanimous dismissal in 1869 once the repealing statute was squarely argued.)

\section{Key Quotations}
\begin{quote}
``It is quite true, as was argued by the counsel for the petitioner, that the appellate jurisdiction of this court is not derived from acts of Congress. It is, strictly speaking, conferred by the Constitution. But it is conferred `with such exceptions and under such regulations as Congress shall make.''' (74 U.S. at 512--513)
\end{quote}

\begin{quote}
``We are not at liberty to inquire into the motives of the legislature. We can only examine into its power under the Constitution, and the power to make exceptions to the appellate jurisdiction of this court is given by express words.'' (74 U.S. at 514)
\end{quote}

\begin{quote}
``Without jurisdiction, the court cannot proceed at all in any cause. Jurisdiction is power to declare the law, and, when it ceases to exist, the only function remaining to the court is that of announcing the fact and dismissing the cause.'' (74 U.S. at 514)
\end{quote}

\begin{quote}
``It is quite clear, therefore, that this court cannot proceed to pronounce judgment in this case, for it has no longer jurisdiction of the appeal, and judicial duty is not less fitly performed by declining ungranted jurisdiction than in exercising firmly that which the Constitution and the laws confer.'' (74 U.S. at 515)
\end{quote}

\begin{quote}
``Counsel seem to have supposed, if effect be given to the repealing act in question, that the whole appellate power of the court, in cases of habeas corpus, is denied. But this is an error. The act of 1868 does not except from that jurisdiction any cases but appeals from Circuit Courts under the act of 1867. It does not affect the jurisdiction which was previously exercised.'' (74 U.S. at 515)
\end{quote}

\section{Doctrinal Significance and Modern Legacy (Through 2026)}
\textit{Ex parte McCardle} is the foundational precedent on congressional ``jurisdiction stripping'' under the Exceptions Clause, yet its actual holding has been cabined by four major doctrinal guardrails developed from 1869 through 2026:
\begin{itemize}
  \item \textbf{The \textit{Yerger} Escape Valve:} Just months after \textit{McCardle}, the Court confirmed its closing hint in \href{https://supreme.justia.com/cases/federal/us/75/85/}{\textit{Ex parte Yerger}}, 75 U.S. (8 Wall.) 85 (1869), holding that the 1868 Act left untouched the Court's appellate authority to review military detention via habeas corpus and certiorari under the Judiciary Act of 1789. More than a century later, in \href{https://supreme.justia.com/cases/federal/us/518/651/}{\textit{Felker v. Turpin}}, 518 U.S. 651 (1996), a unanimous Court invoked \textit{Yerger} to uphold AEDPA's bar on certiorari review of gatekeeper orders denying successive habeas petitions precisely because the statute left intact the Court's authority to entertain ``original'' habeas petitions.
  \item \textbf{The \textit{Klein} Separation-of-Powers Limit:} In \href{https://supreme.justia.com/cases/federal/us/80/128/}{\textit{United States v. Klein}}, 80 U.S. (13 Wall.) 128 (1871), Chief Justice Chase drew a critical line: Congress may withdraw appellate jurisdiction over a class of cases (\textit{McCardle}), but it may not use a purported jurisdictional strip to prescribe a ``rule of decision'' in pending litigation or impair an independent constitutional power (there, the presidential pardon power).
  \item \textbf{External Constitutional Constraints (Suspension Clause \& Due Process):} In \href{https://supreme.justia.com/cases/federal/us/553/723/}{\textit{Boumediene v. Bush}}, 553 U.S. 723 (2008), the Court invalidated Section 7 of the Military Commissions Act of 2006, holding that Congress cannot strip federal habeas jurisdiction over Guantanamo detainees without providing an adequate substitute or complying with the Suspension Clause (Art.\ I, \S~9, cl.~2).
  \item \textbf{Modern Jurisdiction-Stripping and Clear-Statement Doctrine (2018--2026):} In \href{https://supreme.justia.com/cases/federal/us/583/16-498/}{\textit{Patchak v. Zinke}}, 583 U.S. 244 (2018), a plurality led by Justice Thomas reaffirmed \textit{McCardle} and upheld the Gun Lake Act's stripping of federal jurisdiction over pending suits relating to a specific parcel of tribal trust land, reasoning that stripping jurisdiction changes the law rather than exercising judicial power (over Chief Justice Roberts's dissent warning that targeting a single pending suit violates \textit{Klein}). At the same time, the Roberts Court has policed statutory jurisdiction strips through a strict clear-statement rule---as in \href{https://supreme.justia.com/cases/federal/us/598/21-1436/}{\textit{Santos-Zacaria v. Garland}}, 598 U.S. 411 (2023) (treating immigration exhaustion rules as nonjurisdictional absent unmistakable congressional text)---while scholarly and legislative debates through 2026 over Supreme Court reform continue to contest whether Henry Hart's ``essential functions'' thesis bars Congress from using \textit{McCardle} to strip the Court of jurisdiction over entire categories of constitutional rights (see \href{https://texaslawreview.org/court-stripping-court-packing-and-court-defying-revisiting-the-supreme-courts-essential-functions/}{\textit{Texas Law Review}}).
\end{itemize}

\section{Pedagogical Critiques and Exam Hypotheticals}
\subsection{Key Critiques}
\begin{itemize}
  \item \textbf{Broad vs.\ Narrow Reading:} Read broadly, \textit{McCardle}'s text (``power to make exceptions \dots is given by express words'') suggests plenary congressional authority to eliminate Supreme Court appellate review over controversial constitutional topics (e.g., school prayer, abortion, or gun regulations). Read narrowly, \textit{McCardle} stands only for the modest proposition that Congress may eliminate \emph{one} statutory avenue of appellate review when another avenue (here, the 1789 Act) remains open.
  \item \textbf{The ``Essential Functions'' Thesis:} Professor Henry Hart famously argued that ``exceptions'' cannot be allowed to ``engulf the rule'' or destroy the Supreme Court's essential role in ensuring the supremacy and uniformity of federal law (\textit{The Power of Congress to Limit the Jurisdiction of Federal Courts}, 66 Harv.\ L.\ Rev.\ 1362, 1364--65 (1953)). Critics of Hart counter that the text of Article III, Section 2 contains no ``essential functions'' carve-out and that the Madisonian Compromise left lower federal courts entirely to congressional discretion.
\end{itemize}

\subsection{Exam Hypothetical}
Congress enacts the ``Uniform Flag Protection Act,'' making flag burning a federal crime and adding Section 5: ``Pursuant to Article III, Section 2, the Supreme Court shall have no appellate jurisdiction, by certiorari, appeal, or original writ of habeas corpus, to review any federal or state court decision upholding a conviction under this Act, including cases currently pending.'' A defendant convicted in state court appeals to the U.S.\ Supreme Court, which already granted certiorari before Section 5 passed.
\begin{itemize}
  \item \textbf{Analysis:} Unlike \textit{McCardle} and \textit{Felker}, Section 5 forecloses \emph{all} Supreme Court review (eliminating the \textit{Yerger} escape valve) and discriminates asymmetrically by stripping jurisdiction only over decisions \emph{upholding} convictions (raising \textit{Klein} rule-of-decision, Due Process, and First Amendment viewpoint-discrimination objections, as well as a Suspension Clause challenge to stripping habeas). Even if a court rejects the ``essential functions'' thesis under \textit{McCardle}'s plain-text rationale, Section 5 is vulnerable under external constitutional limits (First and Fifth Amendments and the Suspension Clause) and \textit{Klein}'s prohibition on Congress rigging outcomes while leaving courts open to decide in the government's favor.
  \item \textbf{Opposite Side Contrast (\textit{United States v. Klein}, 80 U.S. (13 Wall.) 128 (1871)):} Contrast \textit{McCardle} with \textit{Klein}, decided just two years later by the same Chief Justice. In \textit{McCardle}, Congress neutrally withdrew one statutory route of appellate jurisdiction without telling the Court how to evaluate the merits; the Court dismissed. In \textit{Klein}, Congress left courts open to adjudicate claims but ordered them to treat a presidential pardon as conclusive proof of disloyalty and then dismiss on jurisdictional grounds; the Court struck the statute down. On an exam, always ask: Is Congress genuinely closing the courthouse door (\textit{McCardle}/\textit{Patchak}), or is it keeping the door open while dictating the outcome or violating an independent constitutional right (\textit{Klein}/\textit{Boumediene})?
\end{itemize}

\end{document}
